% !PRINTED 160
\weberpage{160}

%<<CONT
\[
1 - \varkappa^2\operatorname{sn}^2(u\pm a)\operatorname{sn}^2(v\pm a) = \dfrac{\left[1-\varkappa^2\operatorname{sn}^4\left(\dfrac{u+v}{2}\pm a\right)\right]\left(1-\varkappa^2\operatorname{sn}^4\dfrac{u-v}{2}\right)}{\left[1-\varkappa^2\operatorname{sn}^2\left(\dfrac{u+v}{2}\pm a\right)\operatorname{sn}^2\!\left(\dfrac{u-v}{2}\right)\right]^2},
\]
woraus sich die Übereinstimmung von (9) mit (7) ergibt.

\begin{center}
\textbf{\S\,49. \gesperrt{Die Transzendenten zweiter und dritter Gattung}\\
\gesperrt{von Weierstrass.}}
\end{center}

Es sind nun noch die Transzendenten zweiter und dritter
Gattung in der \gesperrt{Weierstrassschen} Form aufzustellen. Zu der
ersteren gelangen wir ähnlich wie in \S\,47 durch zweimalige
Differentiation von $\log \sigma(u)$. Es ergibt sich so aus
\weq{1}{\sigma(u) = \omega_1 e^{\frac{\eta_1 u^2}{\omega_1}} \dfrac{\vartheta_{11}\left(\dfrac{u}{\omega_1}\right)}{\vartheta_{11}'} \ [\S\,37,\,(4)],}
\[
\begin{aligned}
\frac{d^2 \log \sigma(u)}{du^2} &= \frac{2\eta_1}{\omega_1} + \frac{d^2\log \vartheta_{11}\left(\dfrac{u}{\omega_1}\right)}{du^2},\\
&= \frac{2\eta_1}{\omega_1} + \frac{1}{\omega_1^2}\frac{\vartheta_{00}''}{\vartheta_{00}} - \frac{1}{\omega_1^2}\frac{\vartheta_{11}'^2}{\vartheta_{00}^2}\frac{\vartheta_{00}^2\left(\dfrac{u}{\omega_1}\right)}{\vartheta_{11}^2\left(\dfrac{u}{\omega_1}\right)} \ [\S\,23,\,(18)]
\end{aligned}
\]
und also, wenn man den $\vartheta$-Quotienten nach \S\,42 durch elliptische
Funktionen und diese nach \S\,46 durch die $\wp$-Funktion ausdrückt:
\[
\frac{d^2\log\sigma(u)}{du^2} = A - \wp(u),
\]
worin $A$ eine Konstante ist. Diese findet man aber, wenn man
die Gleichung nach \S\,46, (8) so schreibt:
\[
\sigma''(u)\sigma(u) - \sigma'(u)\sigma'(u) = A\sigma^2(u) - \tfrac{1}{3}\left[\sigma_{00}^2(u) + \sigma_{01}^2(u) + \sigma_{10}^2(u)\right].
\]
Wenn man hierin zweimal differentiiert und dann $u=0$
setzt, so ergibt sich $A=0$ [\S\,35, (12), \S\,37, (5), (7)], und wir
erhalten:
\weq{2}{\frac{d^2\log\sigma(u)}{du^2} = -\wp(u),}
eine Formel, die bei \gesperrt{Weierstrass} zur Definition der $\wp$-Funktion
dient.
